% subsubsection: 1次分数型  \, $\displaystyle a_{n+1}=\frac{pa_n+q}{ra_n+s}$

\subsubsection{\texorpdfstring{一次分数型  \, $\displaystyle a_{n+1}=\frac{pa_n+q}{ra_n+s}$}{一次分数型}}\label{節：1次分数型}
    分子にも定数項がある\bookterm{recurrence}{漸化式}を考えよう.
    以下では$r\ne0$, $ps-qr\ne0$とし,元の分母が$0$にならない範囲で議論する.\footnote{$r=0,\ s\ne0$なら定数項を消す型へ戻る.
    $r\ne0,\ ps-qr=0$なら,分母が非零の入力に対して右辺は定数$p/r$になる.次の分母も確認する.例えば$(a_n-1)/(a_n-1)$は,初項が$1$でなくても次項が$1$になり,次の更新は定義できない.}
    分子の定数項$q$を消せれば,逆数型として解ける.

    \bookstep{一つの差を追う}
    \sectionref{節：特性方程式型}と同じく,各項から定数を引いてみよう.
    $u_n=a_n-\alpha$と置き,分子の定数項が消えるように$\alpha$を選ぶ.
    $a_n=u_n+\alpha$を代入すると,
    \BookMathLayout{\begin{align*}
      u_{n+1}&=\frac{p(u_n+\alpha)+q}{r(u_n+\alpha)+s}-\alpha\\
      &=\frac{(p-r\alpha)u_n+p\alpha+q-r\alpha^2-s\alpha}
        {ru_n+r\alpha+s}
    \end{align*}}{\[
      \begin{aligned}
        &u_{n+1}\\
        &{}=\frac{p(u_n+\alpha)+q}{r(u_n+\alpha)+s}-\alpha\\
        &{}=\frac{(p-r\alpha)u_n+p\alpha+q-r\alpha^2-s\alpha}
          {ru_n+r\alpha+s}
      \end{aligned}
    \]}
    従って,
    \[
      r\alpha^2+(s-p)\alpha-q=0
      \quad\Longleftrightarrow\quad
      \alpha=\frac{p\alpha+q}{r\alpha+s}
    \]
    を満たす$\alpha$,すなわち元の\bookterm{fixed}{不動点}を引けば,
    \[
      u_{n+1}=\frac{(p-r\alpha)u_n}{ru_n+r\alpha+s}
    \]
    という逆数型になる.ここで,$p-r\alpha$も$r\alpha+s$も零でない.\footnote{不動点の方程式より,$(p-r\alpha)(r\alpha+s)=ps-qr\ne0$.}
    初項が$\alpha$なら定数列である.それ以外では$u_n$は途中で零にならず,逆数を取って解ける.

    \bookstep{二つの差の比を追う}
    不動点の方程式
    \[
      rx^2+(s-p)x-q=0
    \]
    が異なる二つの実数解$\alpha,\beta$を持つとしよう.
    どちらも定数項を消せるので,先ほどの変形を$a_n$で書けば,
    \begin{align*}
      a_{n+1}-\alpha
      &=\frac{p-r\alpha}{ra_n+s}(a_n-\alpha),\\
      a_{n+1}-\beta
      &=\frac{p-r\beta}{ra_n+s}(a_n-\beta)
    \end{align*}
    となる.倍率に$a_n$が含まれるため,各々の差は一般には\bookterm{geometric}{等比数列}にならない.
    しかし,\textbf{二つの式の比を取れば,共通の分母が消える}.

    初項が$\alpha$または$\beta$なら定数\bookterm{sequence}{数列}として分けておく.
    それ以外では,$p-r\beta\ne0$なので,$a_n\ne\beta$から$a_{n+1}\ne\beta$が続く.
    従って,差の比を取れる.
    \begin{align*}
      \frac{a_{n+1}-\alpha}{a_{n+1}-\beta}
      &=\frac{p-r\alpha}{p-r\beta}
        \frac{a_n-\alpha}{a_n-\beta}
    \end{align*}
    ここで$b_n=(a_n-\alpha)/(a_n-\beta)$と置けば,
    \[
      b_{n+1}=\frac{p-r\alpha}{p-r\beta}b_n
    \]
    となる.一つの差を逆数型として解く代わりに,二つの差の比を直接等比型にできた.\footnote{二つの解のどちらを$\alpha,\beta$と名付けてもよい.}
    初項$b_1=(a_1-\alpha)/(a_1-\beta)$を用いて,
    \[
      b_n=\frac{a_1-\alpha}{a_1-\beta}\left(\frac{p-r\alpha}{p-r\beta}\right)^{n-1} \booktag{mobius-geometric}{＊}
    \]
    を得る.

    \bookstep{元の項へ戻す}
    $b_n(a_n-\beta)=a_n-\alpha$を$a_n$について解くと,
    \[
      a_n=\frac{\beta b_n-\alpha}{b_n-1}
    \]
    となる.有限の$a_n$では$\alpha\ne\beta$より$b_n\ne1$なので,この割り算ができる.
    \bookeqbackref{mobius-geometric}{(＊)}を代入すればよい.
    逆に,公式上で$b_n=1$となる項は復元できず,元の更新をそこまで定義できない.

    不動点の方程式が重解$\alpha$を持つ場合は,異なる差を作れない.
    一つの差$u_n=a_n-\alpha$に戻り,逆数を取れば等差型になる.
    例題(3)で計算し,例題の後で一般の係数でも確かめよう.

    \noindent \bookblocklink{example-mobius-1}{problem}{answer}{【例題】}\par\nobreak
    \begin{enumerate}[label=(\arabic*),parsep=3pt,format=\bookitemlink{example-mobius-1}{problem}{answer}]
      \item $\displaystyle a_1=3,\, a_{n+1}=\frac{2a_n+1}{a_n+2}$
      \item $\displaystyle a_1=0,\, a_{n+1}=\frac{a_n+4}{a_n+1}$
      \item $\displaystyle a_1=5,\, a_{n+1}=\dfrac{4a_{n}-9}{a_{n}-2}$
      \item $\displaystyle a_1=1,\, a_{n+1}=1+\frac{1}{a_n}$
    \end{enumerate}

    \noindent\bookblocklink{example-mobius-1}{hint}{problem}{【方針】}\par\nobreak
    \begin{enumerate}[label=(\arabic*),parsep=3pt,format=\bookitemlink{example-mobius-1}{hint}{problem}]
      \item 不動点$1,-1$からの差の比を作ると,共通の分母$a_n+2$が消えます.
      \item 同様に\bookterm{fixed}{不動点}$\alpha,\beta$を求めましょう.今回は片方が負の値になります.
      \item 今回は\bookterm{fixed}{不動点}の方程式が重解になります.\,$\dfrac{a_n-\alpha}{a_n-\beta}$は使えないので,\,$b_n=a_n-\alpha$と置いてから逆数型に帰着させる別の手法を使います.
      \item このような形も一次分数型です.解と係数の関係を活用して解きましょう.
    \end{enumerate}

    \noindent\bookblocklink{example-mobius-1}{answer}{problem}{【解答】}\par\nobreak
    \begin{enumerate}[label=(\arabic*),parsep=3pt,format=\bookitemlink{example-mobius-1}{answer}{problem}]
      \item $\displaystyle a_n=\frac{2\cdot3^{n-1}+1}{2\cdot3^{n-1}-1}$
      \item $\displaystyle a_n=\frac{2\left(3^{n-1}+(-1)^n\right)}{3^{n-1}-(-1)^n}$
      \item $a_{n}=\dfrac{6n-1}{2n-1}$
      \item $\displaystyle a_n=\cfrac{\left(\cfrac{1-\sqrt{5}}{2}\right)^{n+1}-\left(\cfrac{1+\sqrt{5}}{2}\right)^{n+1}}{\left(\cfrac{1-\sqrt{5}}{2}\right)^n-\left(\cfrac{1+\sqrt{5}}{2}\right)^n}$
    \end{enumerate}

    \noindent\bookblocklink{example-mobius-1}{solution}{problem}{【解法】}\par\nobreak
    \begin{enumerate}[label=(\arabic*),parsep=3pt,format=\bookitemlink{example-mobius-1}{solution}{problem}]
      \item $\displaystyle a_1=3,\, a_{n+1}=\frac{2a_n+1}{a_n+2}$\\
            \bookterm{fixed}{不動点}の方程式$x=(2x+1)/(x+2)$の解は$1,-1$である.
            $b_n=(a_n-1)/(a_n+1)$と置くと,
            \begin{align*}
              b_{n+1}&=\cfrac{\cfrac{2a_n+1}{a_n+2}-1}{\cfrac{2a_n+1}{a_n+2}+1}\\
              &=\frac{2a_n+1-(a_n+2)}{2a_n+1+(a_n+2)}\\
              &=\frac13\frac{a_n-1}{a_n+1}=\frac13b_n
            \end{align*}
            $b_1=(3-1)/(3+1)=1/2$なので,
            \[
              b_n=\frac12\left(\frac13\right)^{n-1}
                =\frac1{2\cdot3^{n-1}}
            \]
            である.元の項への復元は,
            \begin{align*}
              b_n(a_n+1)&=a_n-1,\\
              a_n&=\frac{1+b_n}{1-b_n}
                =\frac{2\cdot3^{n-1}+1}{2\cdot3^{n-1}-1}
            \end{align*}
            となる.
      \item $\displaystyle a_1=0,\, a_{n+1}=\frac{a_n+4}{a_n+1}$\\
            \bookterm{fixed}{不動点}の方程式$x=(x+4)/(x+1)$の解は$2,-2$である.
            倍率は$(1-2)/(1+2)=-1/3$なので,$b_n=(a_n-2)/(a_n+2)$と置くと,
            \begin{align*}
              b_{n+1}&=-\frac13b_n,\quad b_1=-1,\\
              b_n&=(-1)\left(-\frac13\right)^{n-1}
                =\frac{(-1)^n}{3^{n-1}}
            \end{align*}
            を得る.$b_n(a_n+2)=a_n-2$から復元して,
            \begin{align*}
              a_n&=\frac{2(1+b_n)}{1-b_n}\\
              &=\frac{2\left(3^{n-1}+(-1)^n\right)}{3^{n-1}-(-1)^n}
            \end{align*}
            となる.
      \item $\displaystyle a_1=5,\, a_{n+1}=\dfrac{4a_n-9}{a_n-2}$\\
            \bookterm{fixed}{不動点}の方程式$x=(4x-9)/(x-2)$は重解$3$を持つ.
            $b_n=a_n-3$と置くと,
            \begin{align*}
              b_{n+1}&=\frac{4a_n-9}{a_n-2}-3
                =\frac{a_n-3}{a_n-2}\\
              &=\frac{b_n}{b_n+1},\quad b_1=2
            \end{align*}
            となる.正の$b_n$が続くので,$b_n$も元の分母$b_n+1$も零でない.
            逆数を取れば,
            \begin{align*}
              \frac1{b_{n+1}}&=\frac1{b_n}+1,\\
              \frac1{b_n}&=\frac12+(n-1)=\frac{2n-1}{2}
            \end{align*}
            である.従って,
            \[
              a_n=3+b_n=3+\frac2{2n-1}=\frac{6n-1}{2n-1}
            \]
            を得る.
      \item $\displaystyle a_1=1,\, a_{n+1}=1+\frac1{a_n}$\\
            右辺を$(a_n+1)/a_n$と書けば,一次分数型である.
            \bookterm{fixed}{不動点}の方程式
            \[
              x=\frac{x+1}{x}\booktag{mobius-example-fixed}{1}
            \]
            の解を
            \[
              \alpha=\frac{1+\sqrt5}{2},\qquad
              \beta=\frac{1-\sqrt5}{2}
            \]
            とする.\bookeqbackref{mobius-example-fixed}{(1)}の解と係数の関係から,
            \[
              \alpha+\beta=1,\qquad \alpha\beta=-1
                \booktag{mobius-example-roots}{2}
            \]
            である.$b_n=(a_n-\alpha)/(a_n-\beta)$と置くと,
            \begin{align*}
              b_{n+1}&=\frac{a_n+1-\alpha a_n}{a_n+1-\beta a_n}\\
              &=\frac{(1-\alpha)a_n+1}{(1-\beta)a_n+1}
                \booktag{mobius-example-ratio}{3}
            \end{align*}
            となる.\bookeqbackref{mobius-example-roots}{(2)}を\bookeqbackref{mobius-example-ratio}{(3)}に使えば,
            \begin{align*}
              b_{n+1}&=\frac{\beta a_n-\alpha\beta}{\alpha a_n-\alpha\beta}\\
              &=\frac{\beta(a_n-\alpha)}{\alpha(a_n-\beta)}
                =\frac\beta\alpha b_n
            \end{align*}
            である.$b_1=(1-\alpha)/(1-\beta)=\beta/\alpha$より,
            \[
              b_n=\left(\frac\beta\alpha\right)^n
            \]
            となる.復元式$a_n=(\beta b_n-\alpha)/(b_n-1)$に代入し,分母分子に$\alpha^n$を掛けると,
            \begin{align*}
              a_n&=\frac{\beta^{n+1}-\alpha^{n+1}}{\beta^n-\alpha^n}\\
              &=\cfrac{\left(\cfrac{1-\sqrt5}{2}\right)^{n+1}-\left(\cfrac{1+\sqrt5}{2}\right)^{n+1}}
                {\left(\cfrac{1-\sqrt5}{2}\right)^n-\left(\cfrac{1+\sqrt5}{2}\right)^n}
            \end{align*}
            を得る.
    \end{enumerate}
    \bookstep{一つの不動点から進む}
    不動点の方程式が重解$\alpha$を持つときは,差の比を作れない.
    $\alpha=(p-s)/(2r)$なので,$c=p-r\alpha=r\alpha+s\ne0$と置けば,
    \[
      a_{n+1}-\alpha
      =\frac{c(a_n-\alpha)}{r(a_n-\alpha)+c}
    \]
    となる.$a_1=\alpha$なら定数\bookterm{sequence}{数列}として分ける.
    それ以外では,分母が非零の範囲で$c\ne0$より
    \[
      a_{n+1}-\alpha=0
      \quad\Longleftrightarrow\quad a_n-\alpha=0
    \]
    であり,途中で$\alpha$に到達しない.従って逆数を取れて,
    \[
      \frac1{a_{n+1}-\alpha}
      =\frac1{a_n-\alpha}+\frac rc
    \]
    という\bookterm{arithmetic}{等差数列}に帰着する.
    \textbf{異なる不動点なら二つの差の比,重解なら一つの差の逆数}を追えばよい.
    例題(3)は$\alpha=3,\ r=c=1$なので,逆数が毎回$1$ずつ増えた.

    \bookstep{比から見える数列}
    例題(4)は,フィボナッチ\bookterm{sequence}{数列}$F_1=F_2=1,\ F_{n+2}=F_{n+1}+F_n$とつながる.
    隣接項の比$F_{n+1}/F_n$は同じ更新$a_{n+1}=1+1/a_n$と初項$1$を満たすので,$a_n=F_{n+1}/F_n$である.
    この比の\bookterm{limit}{極限}は黄金比$(1+\sqrt5)/2$となる.\footnote{リュカ数列も同じ\bookterm{characteristic}{特性方程式}$x^2=x+1$を持ち,隣接項の比が黄金比へ収束する.
    トリボナッチ数列では$x^3=x^2+x+1$となり,比の極限は約$1.839$のトリボナッチ定数になる.}

  \bookstep{項を並べて予想する}
  一次分数型では,初めの数項から一般項を予想し,帰納法で確かめる方法もある.
  数項を求める誘導があれば,分母と分子の増え方にも注目しよう.例題(3)で試す.\\
    \noindent\bookblocklink{example-mobius-2}{problem}{answer}{【例題】}\par\nobreak
    \begin{quote}
      以下の\bookterm{recurrence}{漸化式}
      \[
        a_1=5,\, a_{n+1}=\dfrac{4a_{n}-9}{a_{n}-2}
      \]
      を満たす\bookterm{sequence}{数列}$a_n$の\bookterm{general}{一般項}を帰納法を用いて求めよ.
    \end{quote}
    \bookblocklink{example-mobius-2}{hint}{problem}{【方針】}
    \begin{quote}
      実験的に代入計算をしてみると,
      \[
        a_1=\frac{5}{1},\,a_2=\frac{11}{3},\,a_3=\frac{17}{5},\,a_4=\frac{23}{7},\,a_5=\frac{29}{9}
      \]
      となる.分母が2ずつ,分子が6ずつ増えていることから,\bookterm{general}{一般項}は
      \[
        a_n=\frac{6n-1}{2n-1}
      \]
      と予想する.
    \end{quote}
    \bookblocklink{example-mobius-2}{answer}{problem}{【解答】}
    \begin{quote}
      予想した\bookterm{general}{一般項}を確かめる.$n=1$では$(6-1)/(2-1)=5$なので成立する.
      正の整数$k$について,
      \[
        a_k=\frac{6k-1}{2k-1}
      \]
      が成り立つと仮定する.このとき,\bookterm{recurrence}{漸化式}より,
      \begin{align*}
        a_{k+1}&=\frac{4(6k-1)-9(2k-1)}{6k-1-2(2k-1)}\\
        &=\frac{6k+5}{2k+1}
          =\frac{6(k+1)-1}{2(k+1)-1}
      \end{align*}
      分母$a_k-2=(2k+1)/(2k-1)$も零でなく,$n=k+1$で成立する.
      帰納法より,$a_n=(6n-1)/(2n-1)$がすべての正の整数$n$で成り立つ.\qed

  \end{quote}
